% Imporditud: tools/import_eksperimendid.py
% Allikas: Eesti füüsikaolümpiaadi eksperimendi ülesannete
%   kogu (Rael Kalda, 2023), ülesanne Ü82, lahendus L82.
\setAuthor{EFO žürii}
\setRound{lõppvoor}
\setYear{2000}
\setNumber{P E1}
\setDifficulty{3}
\setTopic{geomeetriline optika}
\setVahendid{Klaasnõu, anum veega, pliiats}
\setPunktid{10}
\setAllikas{Eksperimendi kogu Ü82, lahendus lk 68}
\setLahendusseis{toores}

\prob{Pliiats vees}
Asetage pliiats tühja klaasnõusse ja vaadake seda küljelt. Valage nõu poolenisti vett täis. Mida märkate? Kas nähtused olenevad pliiatsi asendist? Kui olenevad, siis kuidas? Põhjendage vastuseid.

\hint

\solu
Vee lisamisel esineb pliiatsi ``murdumine'' --- (1 p) Selle põhjendus --- (1 p) Nähtus oleneb pliiatsi asendist --- (1 p) Suurim on siis, kui pliiats on klaasi keskkohast kaugeimas asukohas --- (1 p)

Põhjendus: jutt --- (1 p), joonis --- (1 p) Nähtus puudub kui pliiats asub vertikaalselt (1 p) klaasi keskel (1 p) Põhjendus: jutt --- (1 p), joonis --- (1 p)
\probend
